\chapter{Experimental Setup}

\section{Datasets and Tasks}

All data are loaded from \texttt{nyu-mll/glue} through the Hugging Face \texttt{datasets} library~\cite{lhoest2021datasets}. Table~\ref{tab:data} lists the five tasks. The official training split of each task is permuted once with a fixed seed and divided into a few-shot pool of 256 labelled examples for router training, a development split of 10\% capped at 1,000 examples, and the remaining expert-training data. The experts never see the pool or the development split. The development split makes every model choice (expert epoch, Task Arithmetic $\lambda$, router step), and the official validation split is used only for reported scores. GLUE test labels are not used.

\begin{table}[htbp]
  \centering\tablestyle
  \caption{GLUE tasks and data splits. Train and validation sizes and majority rates were verified against the dataset; other sizes follow from the split rule.}
  \label{tab:data}
  \begin{tabularx}{\linewidth}{l X l r r r r r r}
    \toprule
    Task & Input & Task score & Train & Pool & Dev & Expert train & Validation & Majority (\%)\\
    \midrule
    MRPC & Sentence pair & mean(acc, F1) & 3,668 & 256 & 367 & 3,045 & 408 & 68.4\\
    RTE & Sentence pair & accuracy & 2,490 & 256 & 249 & 1,985 & 277 & 52.7\\
    CoLA & Sentence & MCC & 8,551 & 256 & 855 & 7,440 & 1,043 & 69.1\\
    SST-2 & Sentence & accuracy & 67,349 & 256 & 1,000 & 66,093 & 872 & 50.9\\
    QNLI & Question, sentence & accuracy & 104,743 & 256 & 1,000 & 103,487 & 5,463 & 50.5\\
    \bottomrule
  \end{tabularx}
\end{table}

\section{Implementation Environment}

The project is implemented in Python using PyTorch, Hugging Face Transformers~\cite{wolf2020transformers}, \texttt{datasets} and scikit-learn. The reported run used Google Colab with one NVIDIA Tesla T4 GPU (15\,GB), with torch 2.11.0, transformers 5.20.2, datasets 4.8.5 and scikit-learn 1.6.1. The project's unit tests pass on CPU, but they use small random models and do not reproduce the experiments.

\section{Experimental Configuration}
\label{sec:config}

Table~\ref{tab:config} gives the configuration. Both stages use mixed precision. Tokenisation truncates to 128 tokens and pads dynamically. The Task Arithmetic grid covers $\lambda\in\{0.1,0.2,\dots,1.0\}$, and $\lambda=0.3$ gave the best development score (62.8). Twelve FSFT runs were performed: full-MLP routers with $k\in\{16,64,256\}$ and intermediate-only routers with $k=256$, each with seeds 0, 1 and 2. Seeds vary the router initialisation and the few-shot sample only.

\begin{table}[htbp]
  \centering\tablestyle
  \caption{Training configuration of the two optimisation stages.}
  \label{tab:config}
  \begin{tabularx}{\linewidth}{>{\bfseries}l *{2}{>{\raggedright\arraybackslash}X}}
    \toprule
    Setting & Task-specific experts & FSFT (routers)\\
    \midrule
    Trainable parameters & All of BERT and the head (109.5M per expert) & Routers only (7.13M, shared across tasks)\\
    Optimiser & AdamW~\cite{loshchilov2019adamw}; weight decay 0.01, none on 1-D parameters & AdamW; no weight decay\\
    Learning rate & $2\times10^{-5}$ & $1\times10^{-4}$\\
    Schedule & Linear decay, 10\% warm-up & Linear decay, 10\% warm-up\\
    Batch size & 32 & 16 per task (one batch from each task per step)\\
    Duration & 3 epochs & 500 steps\\
    Gradient clipping & Norm 1.0 & None\\
    Model selection & Best epoch by task score on dev & Best step by mean score on dev, every 50 steps\\
    \bottomrule
  \end{tabularx}
\end{table}

\section{Evaluation Metrics}

Scores follow GLUE conventions and are reported multiplied by 100. From true and false positives and negatives, accuracy, F1 and the Matthews Correlation Coefficient~\cite{matthews1975mcc} are
\begin{equation}
\mathrm{Acc}=\frac{TP+TN}{TP+TN+FP+FN},\quad
F_1=\frac{2\,TP}{2\,TP+FP+FN},
\end{equation}
\begin{equation}\mathrm{MCC}=\frac{TP\cdot TN-FP\cdot FN}{\sqrt{(TP+FP)(TP+FN)(TN+FP)(TN+FN)}} .
\end{equation}
The task score $S_t$ is the mean of accuracy and F1 for MRPC, MCC for CoLA and accuracy for the other tasks. Across tasks, the \emph{average score} of a model $m$ is the mean of $S_t(m)$, and its \emph{normalised score} is the mean of $S_t(m)/S_t(\mathrm{expert})$. Because a constant prediction already scores about 75 on MRPC and 50 on the accuracy tasks, the code also computes the \emph{gain over majority}:
\begin{equation}
\mathrm{Gain}(m)=\frac{1}{E}\sum_{t=1}^{E}\frac{S_t(m)-S_t(\mathrm{majority})}{S_t(\mathrm{expert})-S_t(\mathrm{majority})},
\label{eq:gain}
\end{equation}
where 0 is a constant guess and 100\% is the expert. Efficiency is measured by stored parameters and size at 32-bit precision.
